\section{Background and Related Work}
\label{sec:background}

\subsection{Force closes}
A channel is a 2-of-2 multisignature output on the chain. Each party holds its own version of the latest commitment
transaction, which spends that output to the two balances and to one output per pending payment (HTLC). In the
holder's version, the holder's own balance (\texttt{to\_local}) is locked for \texttt{to\_self\_delay} blocks, a value
the peer chose at opening, while the peer's balance (\texttt{to\_remote}) is not~\cite{bolt3}. The delay gives the peer
time to react if the holder broadcasts an old, revoked commitment: the peer can then take every output with a penalty
(justice) transaction~\cite{poon2016bitcoin}. A force close is the broadcast of a commitment; the broadcaster waits out
the delay before sweeping its balance.

Pending payments add deadlines. An offered HTLC returns to the sender after its expiry height
(\texttt{cltv\_expiry}) through a pre-signed HTLC-timeout transaction, or, on the peer's commitment, directly; a
received HTLC is claimed with the payment's preimage through an HTLC-success transaction~\cite{bolt3,bolt5}. A node
that forwards a payment holds two HTLCs for it, and the incoming one expires a margin later than the outgoing one,
the \texttt{cltv\_expiry\_delta} the node advertises in its channel updates~\cite{bolt7}. If the node cannot time out
its outgoing HTLC on-chain within that margin, the downstream node can claim it with the preimage while the
upstream one times out the incoming HTLC, and the forwarder pays for a payment it is not paid for. Flood-and-loot,
time-dilation and replacement-cycling attacks all work by making such a transaction confirm too
late~\cite{harris2020flood,riard2020timedilation,riard2023replacement}.

\subsection{Who pays the fee}
Commitment transactions are signed when the channel state changes and broadcast later, at a fee rate nobody knew at
signing. Three formats answer this differently.

\emph{Pre-anchor commitments} carry a fee fixed at signing. The channel's opener proposes it with
\texttt{update\_fee}, and a peer that finds the rate unreasonable can fail the channel~\cite{bolt2}. The broadcaster
cannot raise the fee later, because its own output is time-locked. These disagreements are a commonly cited cause of
force closes when fees move fast.

\emph{Anchor commitments} (specified in 2020~\cite{bolts688anchors}) add two 330-sat anchor outputs, one per party,
which their owner can spend at once with a child transaction that pays for both (child pays for parent, CPFP). Sixteen
blocks after the commitment confirms, anyone can spend either anchor, so that unused anchors do not stay in the UTXO
set~\cite{bolt3}. The peer's \texttt{to\_remote} output carries a one-block delay, so that the anchors are the only
outputs a child can spend at once, and the CPFP carve-out, a Bitcoin Core mempool rule, let each party attach one
child even if the other had filled the package limits~\cite{core31}. Second-stage HTLC transactions are signed so that
the holder can add inputs to pay their fee~\cite{bolt3}. Simple taproot channels follow the same design with taproot
outputs and anchor scripts of their own~\cite{bolts995taproot}.

\emph{Zero-fee commitments}~\cite{bolts1228zerofee} pay no fee at all. The commitment is a version-3 (TRUC)
transaction~\cite{bip431} with one keyless pay-to-anchor (P2A) output in place of the two anchors, and it is relayed
together with a child that pays for both, which Bitcoin Core supports since version 28.0~\cite{core28}. TRUC's topology limits replace
the carve-out as the defence against transaction pinning~\cite{optechpinning}, and Bitcoin Core 31.0 removed the
carve-out~\cite{core31}. The specification merged the format in May~2026; channels move to it as operators upgrade
their software, which in the Lightning Network takes months to years~\cite{kumar2026upgradedynamics}.

\subsection{Related work}
BitMEX Research first counted non-cooperative closes from the chain, private channels included, and found more than
60,000 by January~2020~\cite{bitmex2020closures}. Zabka et al.\ classified how public channels close
\cite{zabka2022empirical}; Grötschla et al.\ followed public channels from funding to the spending of their closing
outputs~\cite{grotschla2025lifecycle}; Antonelli et al.\ predict which channels will close and
how~\cite{antonelli2026closures}. Kappos et al.\ and Romiti et al.\ used closing transactions to link channels to
on-chain wallets~\cite{kappos2021empirical,romiti2021crosslayer}. Our previous paper measured force closes around
the 2026 Core Lightning shutdown~\cite{kumar2026clientdiversity}. None of these accounts for what a force close costs
or follows the anchor outputs.

The fee-bumping design has been argued from small samples. Daftuar analysed the 2023 transactions that spent
anchor-like outputs, to test whether Lightning's commitment spends would fit version-3 relay
rules~\cite{daftuar2024v3}. An open LND issue proposes to stop sweeping ``dangling'' anchors, on the argument that
third parties sweep them more efficiently, and gives the break-even fee rate without
data~\cite{lnd8997}. The security literature models what happens when closing transactions confirm late: mass closes
that congest the chain~\cite{harris2020flood,sguanci2023massexit}, nodes kept from seeing new
blocks~\cite{riard2020timedilation}, pinned or cycled
replacements~\cite{optechpinning,riard2023replacement}, HTLC slots jammed by attackers~\cite{mizrahi2021congestion}
and fees taken from revoked commitments~\cite{harding2024irrevocable}. We measure how often, and in what fee
conditions, such deadlines are missed in practice.
